% --- Archon physics formalization source begin ---
% archon:physics
% archon:covers PhyXMiniProblems/problem_phyx_mini_0293.lean
% archon:source-report reports/phyx_mini/problem_phyx_mini_0293.source.json

\chapter{Physics problem phyx\_mini\_0293}
\label{ch:PhyXMiniProblems_problem_phyx_mini_0293}

\paragraph{Problem source.}
\#\# Physical scenario

A sinusoidal wave is sent along a string with a linear density of \$2.0\$ g/m. As it travels, the kinetic energies of the mass elements along the string vary. Figure gives the rate \$dK/dt\$ at which kinetic energy passes through the string elements at a particular instant, plotted as a function of distance \$x\$ along the string. Figure also gives the rate at which kinetic energy passes through a particular mass element (at a particular location), plotted as a function of time \$t\$. For both figures, the scale on the vertical (rate) axis is set by \$R\_s = 10\$ W.

\#\# Question

What is the amplitude of the wave?

\#\# Answer choices

- A: A: 0.0025 m

- B: B: 0.0028 m

- C: C: 0.0030 m

- D: D: 0.0032 m

\#\# Figure caption (auxiliary; use the image as primary evidence)

The image contains two line graphs placed side by side.

\#\#\# Left Graph:
- **Y-axis Label:** \textbackslash{}( \textbackslash{}frac\{dK\}\{dt\} \textbackslash{}) (W), representing the rate of change of kinetic energy in watts.
- **X-axis Label:** \textbackslash{}( x \textbackslash{}) (m), indicating distance in meters.
- **Graph Line:** A sinusoidal wave pattern that starts at a peak, descends to a trough, and returns to a peak within the range from 0 to 0.2 meters.
- **Additional Label:** \textbackslash{}( R\_s \textbackslash{}) located at the top left corner of the graph.

\#\#\# Right Graph:
- **Y-axis Label:** \textbackslash{}( \textbackslash{}frac\{dK\}\{dt\} \textbackslash{}) (W), the same as the left graph, indicating the rate of change of kinetic energy in watts.
- **X-axis Label:** \textbackslash{}( t \textbackslash{}) (ms), representing time in milliseconds.
- **Graph Line:** Another sinusoidal wave pattern similar to the left graph, covering the range from 0 to 2 milliseconds.
- **Additional Label:** \textbackslash{}( R\_s \textbackslash{}) at the top left corner of the graph.

Both graphs feature thick, dark lines, indicating the primary data being presented. The sinusoidal pattern suggests a periodic relationship in both spatial and temporal domains.

\#\# Dataset metadata

Wave Properties; Physical Model Grounding Reasoning, Spatial Relation Reasoning

\paragraph{Recorded answer/context.}
D: D: 0.0032 m

\paragraph{Figure/image path.}
/root/proposal\_for\_physic/science-mango/phyx\_mini\_run/phyx\_data/test\_image/293.png

\paragraph{Formalization target.}
create a compiling Lean file with sorry bodies at `PhyXMiniProblems/problem\_phyx\_mini\_0293.lean`. The Lean declarations must preserve the physical quantities, dimensions or dimensional roles, figure labels, governing-law hypotheses, and final relation expressed by this problem.
Use Mathlib/Physlib names found through LeanExplore where available. If a physics API is missing, introduce faithful local abstractions rather than scalar placeholder aliases.

\begin{theorem}[Physics formalization target]
\label{thm:physics:phyx_mini_0293:target}
\lean{PhyXMiniProblems.ProblemPhyXMini0293.problem_phyx_mini_0293}
\uses{def:physics:phyx-mini-0293:phyxminiproblems-problemphyxmini0293-lengthinmeters, def:physics:phyx-mini-0293:phyxminiproblems-problemphyxmini0293-sinusoidalstringwavesetup, def:physics:phyx-mini-0293:phyxminiproblems-problemphyxmini0293-matchessinusoidalstringwavescenario, def:physics:phyx-mini-0293:phyxminiproblems-problemphyxmini0293-hasphysicalstringwaveparameters, def:physics:phyx-mini-0293:phyxminiproblems-problemphyxmini0293-satisfiessinusoidalstringwavelaws, def:physics:phyx-mini-0293:phyxminiproblems-problemphyxmini0293-matcheskineticenergyratefigure, def:physics:phyx-mini-0293:phyxminiproblems-problemphyxmini0293-recordeddatasetanswer, def:physics:phyx-mini-0293:phyxminiproblems-problemphyxmini0293-isuniquematchinganswerchoice}
The assigned autoformalize agent should translate this physics problem into Lean declarations in the covered file, with theorem and lemma proof bodies written as `by sorry`.
\end{theorem}
\begin{proof}
This is an autoformalization task, not a proof task. Produce statements that can later be proved by the physics prover without weakening the physical model.
\end{proof}

% --- Archon declaration topology begin ---
\section{Declaration topology}
\begin{definition}[Length Quantity]
  \label{def:physics:phyx-mini-0293:phyxminiproblems-problemphyxmini0293-lengthquantity}
  \lean{PhyXMiniProblems.ProblemPhyXMini0293.LengthQuantity}
  A nonnegative physical length, used for amplitude and wavelength
\end{definition}
\begin{proof}
  This is the declared model definition of Length Quantity.
\end{proof}
\begin{definition}[Time Quantity]
  \label{def:physics:phyx-mini-0293:phyxminiproblems-problemphyxmini0293-timequantity}
  \lean{PhyXMiniProblems.ProblemPhyXMini0293.TimeQuantity}
  A nonnegative physical duration, used for the wave period
\end{definition}
\begin{proof}
  This is the declared model definition of Time Quantity.
\end{proof}
\begin{definition}[Linear Mass Density Quantity]
  \label{def:physics:phyx-mini-0293:phyxminiproblems-problemphyxmini0293-linearmassdensityquantity}
  \lean{PhyXMiniProblems.ProblemPhyXMini0293.LinearMassDensityQuantity}
  A nonnegative linear mass density, with SI unit kilogram per metre
\end{definition}
\begin{proof}
  This is the declared model definition of Linear Mass Density Quantity.
\end{proof}
\begin{definition}[Speed Quantity]
  \label{def:physics:phyx-mini-0293:phyxminiproblems-problemphyxmini0293-speedquantity}
  \lean{PhyXMiniProblems.ProblemPhyXMini0293.SpeedQuantity}
  A nonnegative propagation speed, with SI unit metre per second
\end{definition}
\begin{proof}
  This is the declared model definition of Speed Quantity.
\end{proof}
\begin{definition}[Angular Frequency Quantity]
  \label{def:physics:phyx-mini-0293:phyxminiproblems-problemphyxmini0293-angularfrequencyquantity}
  \lean{PhyXMiniProblems.ProblemPhyXMini0293.AngularFrequencyQuantity}
  A nonnegative angular frequency; radians are dimensionless
\end{definition}
\begin{proof}
  This is the declared model definition of Angular Frequency Quantity.
\end{proof}
\begin{definition}[Power Quantity]
  \label{def:physics:phyx-mini-0293:phyxminiproblems-problemphyxmini0293-powerquantity}
  \lean{PhyXMiniProblems.ProblemPhyXMini0293.PowerQuantity}
  The plotted rate dK/dt is a power. Its dimension is mass * length\textasciicircum{}2 / time\textasciicircum{}3, whose coherent SI unit is the watt
\end{definition}
\begin{proof}
  This is the declared model definition of Power Quantity.
\end{proof}
\begin{definition}[length Readout]
  \label{def:physics:phyx-mini-0293:phyxminiproblems-problemphyxmini0293-lengthreadout}
  \lean{PhyXMiniProblems.ProblemPhyXMini0293.lengthReadout}
  \uses{def:physics:phyx-mini-0293:phyxminiproblems-problemphyxmini0293-lengthquantity}
  Read a physical length as a real scalar in the selected length unit
\end{definition}
\begin{proof}
  This is the declared model definition of length Readout.
\end{proof}
\begin{definition}[time Readout]
  \label{def:physics:phyx-mini-0293:phyxminiproblems-problemphyxmini0293-timereadout}
  \lean{PhyXMiniProblems.ProblemPhyXMini0293.timeReadout}
  \uses{def:physics:phyx-mini-0293:phyxminiproblems-problemphyxmini0293-timequantity}
  Read a duration as a real scalar in the selected time unit
\end{definition}
\begin{proof}
  This is the declared model definition of time Readout.
\end{proof}
\begin{definition}[length In Meters]
  \label{def:physics:phyx-mini-0293:phyxminiproblems-problemphyxmini0293-lengthinmeters}
  \lean{PhyXMiniProblems.ProblemPhyXMini0293.lengthInMeters}
  \uses{def:physics:phyx-mini-0293:phyxminiproblems-problemphyxmini0293-lengthquantity, def:physics:phyx-mini-0293:phyxminiproblems-problemphyxmini0293-lengthreadout}
  Metre readout used for amplitudes, wavelengths, and the left graph
\end{definition}
\begin{proof}
  This is the declared model definition of length In Meters.
\end{proof}
\begin{definition}[time In Seconds]
  \label{def:physics:phyx-mini-0293:phyxminiproblems-problemphyxmini0293-timeinseconds}
  \lean{PhyXMiniProblems.ProblemPhyXMini0293.timeInSeconds}
  \uses{def:physics:phyx-mini-0293:phyxminiproblems-problemphyxmini0293-timequantity, def:physics:phyx-mini-0293:phyxminiproblems-problemphyxmini0293-timereadout}
  Second readout used in the governing wave equations
\end{definition}
\begin{proof}
  This is the declared model definition of time In Seconds.
\end{proof}
\begin{definition}[time In Milliseconds]
  \label{def:physics:phyx-mini-0293:phyxminiproblems-problemphyxmini0293-timeinmilliseconds}
  \lean{PhyXMiniProblems.ProblemPhyXMini0293.timeInMilliseconds}
  \uses{def:physics:phyx-mini-0293:phyxminiproblems-problemphyxmini0293-timequantity, def:physics:phyx-mini-0293:phyxminiproblems-problemphyxmini0293-timereadout}
  Millisecond readout used on the right graph's horizontal axis
\end{definition}
\begin{proof}
  This is the declared model definition of time In Milliseconds.
\end{proof}
\begin{definition}[linear Mass Density In Kilograms Per Meter]
  \label{def:physics:phyx-mini-0293:phyxminiproblems-problemphyxmini0293-linearmassdensityinkilogramspermeter}
  \lean{PhyXMiniProblems.ProblemPhyXMini0293.linearMassDensityInKilogramsPerMeter}
  \uses{def:physics:phyx-mini-0293:phyxminiproblems-problemphyxmini0293-linearmassdensityquantity}
  Kilogram-per-metre readout of the string's linear mass density
\end{definition}
\begin{proof}
  This is the declared model definition of linear Mass Density In Kilograms Per Meter.
\end{proof}
\begin{definition}[speed In Meters Per Second]
  \label{def:physics:phyx-mini-0293:phyxminiproblems-problemphyxmini0293-speedinmeterspersecond}
  \lean{PhyXMiniProblems.ProblemPhyXMini0293.speedInMetersPerSecond}
  \uses{def:physics:phyx-mini-0293:phyxminiproblems-problemphyxmini0293-speedquantity}
  Metre-per-second readout of the wave speed
\end{definition}
\begin{proof}
  This is the declared model definition of speed In Meters Per Second.
\end{proof}
\begin{definition}[angular Frequency In Radians Per Second]
  \label{def:physics:phyx-mini-0293:phyxminiproblems-problemphyxmini0293-angularfrequencyinradianspersecond}
  \lean{PhyXMiniProblems.ProblemPhyXMini0293.angularFrequencyInRadiansPerSecond}
  \uses{def:physics:phyx-mini-0293:phyxminiproblems-problemphyxmini0293-angularfrequencyquantity}
  Radian-per-second readout of angular frequency
\end{definition}
\begin{proof}
  This is the declared model definition of angular Frequency In Radians Per Second.
\end{proof}
\begin{definition}[power In Watts]
  \label{def:physics:phyx-mini-0293:phyxminiproblems-problemphyxmini0293-powerinwatts}
  \lean{PhyXMiniProblems.ProblemPhyXMini0293.powerInWatts}
  \uses{def:physics:phyx-mini-0293:phyxminiproblems-problemphyxmini0293-powerquantity}
  Watt readout of a kinetic-energy transfer rate
\end{definition}
\begin{proof}
  This is the declared model definition of power In Watts.
\end{proof}
\begin{definition}[String Wave Kind]
  \label{def:physics:phyx-mini-0293:phyxminiproblems-problemphyxmini0293-stringwavekind}
  \lean{PhyXMiniProblems.ProblemPhyXMini0293.StringWaveKind}
  The qualitative wave type stated in the problem
\end{definition}
\begin{proof}
  This is the declared model definition of String Wave Kind.
\end{proof}
\begin{definition}[String Model Kind]
  \label{def:physics:phyx-mini-0293:phyxminiproblems-problemphyxmini0293-stringmodelkind}
  \lean{PhyXMiniProblems.ProblemPhyXMini0293.StringModelKind}
  The idealized string model used by the textbook relation
\end{definition}
\begin{proof}
  This is the declared model definition of String Model Kind.
\end{proof}
\begin{definition}[Propagation Direction]
  \label{def:physics:phyx-mini-0293:phyxminiproblems-problemphyxmini0293-propagationdirection}
  \lean{PhyXMiniProblems.ProblemPhyXMini0293.PropagationDirection}
  Either direction of propagation along the string's x axis
\end{definition}
\begin{proof}
  This is the declared model definition of Propagation Direction.
\end{proof}
\begin{definition}[propagation Direction Sign]
  \label{def:physics:phyx-mini-0293:phyxminiproblems-problemphyxmini0293-propagationdirectionsign}
  \lean{PhyXMiniProblems.ProblemPhyXMini0293.propagationDirectionSign}
  \uses{def:physics:phyx-mini-0293:phyxminiproblems-problemphyxmini0293-propagationdirection}
  Sign in the standard traveling-wave phase kx - s * omega t + phi
\end{definition}
\begin{proof}
  This is the declared model definition of propagation Direction Sign.
\end{proof}
\begin{definition}[Vertical Axis Role]
  \label{def:physics:phyx-mini-0293:phyxminiproblems-problemphyxmini0293-verticalaxisrole}
  \lean{PhyXMiniProblems.ProblemPhyXMini0293.VerticalAxisRole}
  The physical rate represented by the common vertical axis
\end{definition}
\begin{proof}
  This is the declared model definition of Vertical Axis Role.
\end{proof}
\begin{definition}[Horizontal Axis Role]
  \label{def:physics:phyx-mini-0293:phyxminiproblems-problemphyxmini0293-horizontalaxisrole}
  \lean{PhyXMiniProblems.ProblemPhyXMini0293.HorizontalAxisRole}
  The two distinct horizontal axes in the supplied figure
\end{definition}
\begin{proof}
  This is the declared model definition of Horizontal Axis Role.
\end{proof}
\begin{definition}[Rate Graph Panel]
  \label{def:physics:phyx-mini-0293:phyxminiproblems-problemphyxmini0293-rategraphpanel}
  \lean{PhyXMiniProblems.ProblemPhyXMini0293.RateGraphPanel}
  The two panels in the supplied figure
\end{definition}
\begin{proof}
  This is the declared model definition of Rate Graph Panel.
\end{proof}
\begin{definition}[Spatial Graph Feature]
  \label{def:physics:phyx-mini-0293:phyxminiproblems-problemphyxmini0293-spatialgraphfeature}
  \lean{PhyXMiniProblems.ProblemPhyXMini0293.SpatialGraphFeature}
  Five successive features visible from x = 0 through x = 0.2 m
\end{definition}
\begin{proof}
  This is the declared model definition of Spatial Graph Feature.
\end{proof}
\begin{definition}[Temporal Graph Feature]
  \label{def:physics:phyx-mini-0293:phyxminiproblems-problemphyxmini0293-temporalgraphfeature}
  \lean{PhyXMiniProblems.ProblemPhyXMini0293.TemporalGraphFeature}
  Five successive features visible from t = 0 through t = 2 ms
\end{definition}
\begin{proof}
  This is the declared model definition of Temporal Graph Feature.
\end{proof}
\begin{definition}[Kinetic Energy Rate Figure]
  \label{def:physics:phyx-mini-0293:phyxminiproblems-problemphyxmini0293-kineticenergyratefigure}
  \lean{PhyXMiniProblems.ProblemPhyXMini0293.KineticEnergyRateFigure}
  \uses{def:physics:phyx-mini-0293:phyxminiproblems-problemphyxmini0293-lengthquantity, def:physics:phyx-mini-0293:phyxminiproblems-problemphyxmini0293-timequantity, def:physics:phyx-mini-0293:phyxminiproblems-problemphyxmini0293-powerquantity, def:physics:phyx-mini-0293:phyxminiproblems-problemphyxmini0293-verticalaxisrole, def:physics:phyx-mini-0293:phyxminiproblems-problemphyxmini0293-horizontalaxisrole, def:physics:phyx-mini-0293:phyxminiproblems-problemphyxmini0293-rategraphpanel, def:physics:phyx-mini-0293:phyxminiproblems-problemphyxmini0293-spatialgraphfeature, def:physics:phyx-mini-0293:phyxminiproblems-problemphyxmini0293-temporalgraphfeature}
  Primary-image data are stored independently from the modeled wave. The left panel's coordinates are physical lengths; the right panel's coordinates are physical durations. Both panels contain physical power values rather than bare dimensionless ordinates. The fixed snapshot time and probe position are retained even though the figure does not print their numerical values
\end{definition}
\begin{proof}
  This is the declared model definition of Kinetic Energy Rate Figure.
\end{proof}
\begin{definition}[Sinusoidal String Wave Setup]
  \label{def:physics:phyx-mini-0293:phyxminiproblems-problemphyxmini0293-sinusoidalstringwavesetup}
  \lean{PhyXMiniProblems.ProblemPhyXMini0293.SinusoidalStringWaveSetup}
  \uses{def:physics:phyx-mini-0293:phyxminiproblems-problemphyxmini0293-lengthquantity, def:physics:phyx-mini-0293:phyxminiproblems-problemphyxmini0293-timequantity, def:physics:phyx-mini-0293:phyxminiproblems-problemphyxmini0293-linearmassdensityquantity, def:physics:phyx-mini-0293:phyxminiproblems-problemphyxmini0293-speedquantity, def:physics:phyx-mini-0293:phyxminiproblems-problemphyxmini0293-angularfrequencyquantity, def:physics:phyx-mini-0293:phyxminiproblems-problemphyxmini0293-powerquantity, def:physics:phyx-mini-0293:phyxminiproblems-problemphyxmini0293-stringwavekind, def:physics:phyx-mini-0293:phyxminiproblems-problemphyxmini0293-stringmodelkind, def:physics:phyx-mini-0293:phyxminiproblems-problemphyxmini0293-propagationdirection, def:physics:phyx-mini-0293:phyxminiproblems-problemphyxmini0293-kineticenergyratefigure}
  The independent physical quantities and observable fields of the string wave. Arguments of transverseDisplacementInMeters and kineticEnergyPassingRate are coherent-SI position-in-metres and time-in-seconds readouts. Neither function is defined from the requested answer
\end{definition}
\begin{proof}
  This is the declared model definition of Sinusoidal String Wave Setup.
\end{proof}
\begin{definition}[Matches Sinusoidal String Wave Scenario]
  \label{def:physics:phyx-mini-0293:phyxminiproblems-problemphyxmini0293-matchessinusoidalstringwavescenario}
  \lean{PhyXMiniProblems.ProblemPhyXMini0293.MatchesSinusoidalStringWaveScenario}
  \uses{def:physics:phyx-mini-0293:phyxminiproblems-problemphyxmini0293-sinusoidalstringwavesetup}
  Qualitative information stated in the physical scenario
\end{definition}
\begin{proof}
  This is the declared model definition of Matches Sinusoidal String Wave Scenario.
\end{proof}
\begin{definition}[Has Physical String Wave Parameters]
  \label{def:physics:phyx-mini-0293:phyxminiproblems-problemphyxmini0293-hasphysicalstringwaveparameters}
  \lean{PhyXMiniProblems.ProblemPhyXMini0293.HasPhysicalStringWaveParameters}
  \uses{def:physics:phyx-mini-0293:phyxminiproblems-problemphyxmini0293-lengthinmeters, def:physics:phyx-mini-0293:phyxminiproblems-problemphyxmini0293-timeinseconds, def:physics:phyx-mini-0293:phyxminiproblems-problemphyxmini0293-linearmassdensityinkilogramspermeter, def:physics:phyx-mini-0293:phyxminiproblems-problemphyxmini0293-speedinmeterspersecond, def:physics:phyx-mini-0293:phyxminiproblems-problemphyxmini0293-angularfrequencyinradianspersecond, def:physics:phyx-mini-0293:phyxminiproblems-problemphyxmini0293-powerinwatts, def:physics:phyx-mini-0293:phyxminiproblems-problemphyxmini0293-sinusoidalstringwavesetup}
  Positivity and nondegeneracy of the independently stored quantities
\end{definition}
\begin{proof}
  This is the declared model definition of Has Physical String Wave Parameters.
\end{proof}
\begin{definition}[Satisfies Sinusoidal String Wave Laws]
  \label{def:physics:phyx-mini-0293:phyxminiproblems-problemphyxmini0293-satisfiessinusoidalstringwavelaws}
  \lean{PhyXMiniProblems.ProblemPhyXMini0293.SatisfiesSinusoidalStringWaveLaws}
  \uses{def:physics:phyx-mini-0293:phyxminiproblems-problemphyxmini0293-lengthinmeters, def:physics:phyx-mini-0293:phyxminiproblems-problemphyxmini0293-timeinseconds, def:physics:phyx-mini-0293:phyxminiproblems-problemphyxmini0293-linearmassdensityinkilogramspermeter, def:physics:phyx-mini-0293:phyxminiproblems-problemphyxmini0293-speedinmeterspersecond, def:physics:phyx-mini-0293:phyxminiproblems-problemphyxmini0293-angularfrequencyinradianspersecond, def:physics:phyx-mini-0293:phyxminiproblems-problemphyxmini0293-powerinwatts, def:physics:phyx-mini-0293:phyxminiproblems-problemphyxmini0293-propagationdirectionsign, def:physics:phyx-mini-0293:phyxminiproblems-problemphyxmini0293-sinusoidalstringwavesetup}
  General laws for a traveling sinusoidal wave on a uniform string: k lambda = 2 pi and omega T = 2 pi; the wave travels one wavelength in one period, v T = lambda; the transverse displacement is sinusoidal; kinetic-energy density is transported at speed v, giving dK/dt = (1/2) mu v omega\textasciicircum{}2 A\textasciicircum{}2 cos\textasciicircum{}2(phase); the corresponding peak rate is (1/2) mu v omega\textasciicircum{}2 A\textasciicircum{}2. These are governing relations for arbitrary parameters. They contain no numerical amplitude or answer-choice value
\end{definition}
\begin{proof}
  This is the declared model definition of Satisfies Sinusoidal String Wave Laws.
\end{proof}
\begin{definition}[Matches Kinetic Energy Rate Figure]
  \label{def:physics:phyx-mini-0293:phyxminiproblems-problemphyxmini0293-matcheskineticenergyratefigure}
  \lean{PhyXMiniProblems.ProblemPhyXMini0293.MatchesKineticEnergyRateFigure}
  \uses{def:physics:phyx-mini-0293:phyxminiproblems-problemphyxmini0293-lengthinmeters, def:physics:phyx-mini-0293:phyxminiproblems-problemphyxmini0293-timeinseconds, def:physics:phyx-mini-0293:phyxminiproblems-problemphyxmini0293-timeinmilliseconds, def:physics:phyx-mini-0293:phyxminiproblems-problemphyxmini0293-linearmassdensityinkilogramspermeter, def:physics:phyx-mini-0293:phyxminiproblems-problemphyxmini0293-powerinwatts, def:physics:phyx-mini-0293:phyxminiproblems-problemphyxmini0293-spatialgraphfeature, def:physics:phyx-mini-0293:phyxminiproblems-problemphyxmini0293-temporalgraphfeature, def:physics:phyx-mini-0293:phyxminiproblems-problemphyxmini0293-sinusoidalstringwavesetup}
  Primary-image evidence and the numerical datum from the problem text. The plotted kinetic-energy rate is proportional to cos\textasciicircum{}2 of the underlying wave phase, so adjacent rate peaks are separated by half a wavelength in the spatial panel and half a wave period in the temporal panel. Those two interpretation equations determine wavelength and period, not amplitude
\end{definition}
\begin{proof}
  This is the declared model definition of Matches Kinetic Energy Rate Figure.
\end{proof}
\begin{lemma}[wavelength and period from rate graphs]
  \label{lem:physics:phyx-mini-0293:phyxminiproblems-problemphyxmini0293-wavelength-and-period-from-rate-graphs}
  \lean{PhyXMiniProblems.ProblemPhyXMini0293.wavelength_and_period_from_rate_graphs}
  \uses{def:physics:phyx-mini-0293:phyxminiproblems-problemphyxmini0293-lengthinmeters, def:physics:phyx-mini-0293:phyxminiproblems-problemphyxmini0293-timeinmilliseconds, def:physics:phyx-mini-0293:phyxminiproblems-problemphyxmini0293-sinusoidalstringwavesetup, def:physics:phyx-mini-0293:phyxminiproblems-problemphyxmini0293-matcheskineticenergyratefigure}
  The two figure periods determine lambda = 0.2 m and T = 2 ms
\end{lemma}
\begin{proof}
  The conclusion follows from the governing relations and figure readouts named in the statement of wavelength and period from rate graphs.
\end{proof}
\begin{lemma}[speed and angular Frequency from rate graphs]
  \label{lem:physics:phyx-mini-0293:phyxminiproblems-problemphyxmini0293-speed-and-angularfrequency-from-rate-graphs}
  \lean{PhyXMiniProblems.ProblemPhyXMini0293.speed_and_angularFrequency_from_rate_graphs}
  \uses{def:physics:phyx-mini-0293:phyxminiproblems-problemphyxmini0293-speedinmeterspersecond, def:physics:phyx-mini-0293:phyxminiproblems-problemphyxmini0293-angularfrequencyinradianspersecond, def:physics:phyx-mini-0293:phyxminiproblems-problemphyxmini0293-sinusoidalstringwavesetup, def:physics:phyx-mini-0293:phyxminiproblems-problemphyxmini0293-hasphysicalstringwaveparameters, def:physics:phyx-mini-0293:phyxminiproblems-problemphyxmini0293-satisfiessinusoidalstringwavelaws, def:physics:phyx-mini-0293:phyxminiproblems-problemphyxmini0293-matcheskineticenergyratefigure}
  The graph-derived scales give v = 100 m/s and omega = 1000 pi rad/s
\end{lemma}
\begin{proof}
  The conclusion follows from the governing relations and figure readouts named in the statement of speed and angular Frequency from rate graphs.
\end{proof}
\begin{definition}[Answer Choice]
  \label{def:physics:phyx-mini-0293:phyxminiproblems-problemphyxmini0293-answerchoice}
  \lean{PhyXMiniProblems.ProblemPhyXMini0293.AnswerChoice}
  Labels attached to the four amplitudes printed with the problem
\end{definition}
\begin{proof}
  This is the declared model definition of Answer Choice.
\end{proof}
\begin{definition}[displayed Amplitude In Meters]
  \label{def:physics:phyx-mini-0293:phyxminiproblems-problemphyxmini0293-displayedamplitudeinmeters}
  \lean{PhyXMiniProblems.ProblemPhyXMini0293.displayedAmplitudeInMeters}
  \uses{def:physics:phyx-mini-0293:phyxminiproblems-problemphyxmini0293-answerchoice}
  Amplitude in metres printed beside each answer label
\end{definition}
\begin{proof}
  This is the declared model definition of displayed Amplitude In Meters.
\end{proof}
\begin{definition}[recorded Dataset Answer]
  \label{def:physics:phyx-mini-0293:phyxminiproblems-problemphyxmini0293-recordeddatasetanswer}
  \lean{PhyXMiniProblems.ProblemPhyXMini0293.recordedDatasetAnswer}
  \uses{def:physics:phyx-mini-0293:phyxminiproblems-problemphyxmini0293-answerchoice}
  Dataset provenance retained separately from the physical premises
\end{definition}
\begin{proof}
  This is the declared model definition of recorded Dataset Answer.
\end{proof}
\begin{definition}[Rounds To Displayed Four Decimal Meters]
  \label{def:physics:phyx-mini-0293:phyxminiproblems-problemphyxmini0293-roundstodisplayedfourdecimalmeters}
  \lean{PhyXMiniProblems.ProblemPhyXMini0293.RoundsToDisplayedFourDecimalMeters}
  \uses{def:physics:phyx-mini-0293:phyxminiproblems-problemphyxmini0293-lengthquantity, def:physics:phyx-mini-0293:phyxminiproblems-problemphyxmini0293-lengthinmeters, def:physics:phyx-mini-0293:phyxminiproblems-problemphyxmini0293-answerchoice, def:physics:phyx-mini-0293:phyxminiproblems-problemphyxmini0293-displayedamplitudeinmeters}
  Agreement after rounding an exact amplitude to four decimal metres
\end{definition}
\begin{proof}
  This is the declared model definition of Rounds To Displayed Four Decimal Meters.
\end{proof}
\begin{definition}[Is Unique Matching Answer Choice]
  \label{def:physics:phyx-mini-0293:phyxminiproblems-problemphyxmini0293-isuniquematchinganswerchoice}
  \lean{PhyXMiniProblems.ProblemPhyXMini0293.IsUniqueMatchingAnswerChoice}
  \uses{def:physics:phyx-mini-0293:phyxminiproblems-problemphyxmini0293-sinusoidalstringwavesetup, def:physics:phyx-mini-0293:phyxminiproblems-problemphyxmini0293-answerchoice, def:physics:phyx-mini-0293:phyxminiproblems-problemphyxmini0293-roundstodisplayedfourdecimalmeters}
  The selected answer is the unique displayed value matching the amplitude
\end{definition}
\begin{proof}
  This is the declared model definition of Is Unique Matching Answer Choice.
\end{proof}
\begin{lemma}[exact amplitude from kinetic Energy rate]
  \label{lem:physics:phyx-mini-0293:phyxminiproblems-problemphyxmini0293-exact-amplitude-from-kineticenergy-rate}
  \lean{PhyXMiniProblems.ProblemPhyXMini0293.exact_amplitude_from_kineticEnergy_rate}
  \uses{def:physics:phyx-mini-0293:phyxminiproblems-problemphyxmini0293-lengthinmeters, def:physics:phyx-mini-0293:phyxminiproblems-problemphyxmini0293-sinusoidalstringwavesetup, def:physics:phyx-mini-0293:phyxminiproblems-problemphyxmini0293-matchessinusoidalstringwavescenario, def:physics:phyx-mini-0293:phyxminiproblems-problemphyxmini0293-hasphysicalstringwaveparameters, def:physics:phyx-mini-0293:phyxminiproblems-problemphyxmini0293-satisfiessinusoidalstringwavelaws, def:physics:phyx-mini-0293:phyxminiproblems-problemphyxmini0293-matcheskineticenergyratefigure}
  From the two graph periods, lambda = 0.2 m, T = 0.002 s, v = 100 m/s, and omega = 1000 pi rad/s. With mu = 0.002 kg/m and peak kinetic-energy passing rate 10 W, the law R\_max = (1/2) mu v omega\textasciicircum{}2 A\textasciicircum{}2 gives the exact amplitude A = 1 / (100 pi) m
\end{lemma}
\begin{proof}
  The conclusion follows from the governing relations and figure readouts named in the statement of exact amplitude from kinetic Energy rate.
\end{proof}
% --- Archon declaration topology end ---
% --- Archon physics formalization source end ---
